\section*{Capítulo \ref{ch:b3:genomics} --- Genómica y secuenciación}

\begin{solution}{exo:b3:genomics:1}
Un \omterm{def:b3:genomics:sanger}{didesoxinucleótido} no tiene hidroxilo 3$'$, de modo que no puede
formarse ningún enlace fosfodiéster con el nucleótido siguiente y la
cadena se detiene. Con solo formas didesoxi toda cadena se detendría en
la primera posición; con solo formas normales no se detendría ninguna. La
mezcla convierte la terminación en un suceso aleatorio en cada posición,
de manera que los productos forman una escalera completa, una longitud
por cada base del molde.
\end{solution}

\begin{solution}{exo:b3:genomics:2}
$4\times 10^{8}\times \qty{300}{bp} = 1.2\times 10^{11}$ bases. Ser
humano: $1.2\times 10^{11}/3.2\times 10^{9} \approx 38\times$. Bacteria:
$1.2\times 10^{11}/5\times 10^{6} = \num{24000}\times$.
\end{solution}

\begin{solution}{exo:b3:genomics:3}
Un contig es una secuencia contigua ensamblada a partir de \omterm{def:b3:genomics:ngs}{lecturas}
solapadas; un \omterm{def:b3:genomics:assembly}{andamio} es un conjunto ordenado y orientado de contigs con
huecos de tamaño estimado entre ellos; el \omterm{def:b3:genomics:assembly}{N50} es la longitud de contig a
la que los contigs ordenados alcanzan la mitad de las bases ensambladas.
Aquí los diez contigs de \qty{8}{Mb} contienen ya \qty{80}{Mb}, más allá
del punto medio de \qty{50}{Mb} (el séptimo llega a \qty{56}{Mb}):
\omterm{def:b3:genomics:assembly}{N50} $= \qty{8}{Mb}$.
\end{solution}

\begin{solution}{exo:b3:genomics:4}
El \omterm{prop:b3:genomics:sizes}{tamaño del genoma} no sigue a la complejidad del organismo: una cebolla
tiene cinco veces el ADN de una persona y un pez pulmonado cuarenta
veces; la levadura y el gusano difieren diez veces en ADN con números de
genes parecidos. El exceso es no codificante: elementos transponibles y
sus restos, intrones y repeticiones satélite.
\end{solution}

\begin{solution}{exo:b3:genomics:5}
$e^{-c} = 10^{-6}$ da $c = 6\ln 10 = 13.8$. Para $c = 8$: $N =
cG/L = 8\times 5\times 10^{6}/150 \approx \num{267000}$ \omterm{def:b3:genomics:ngs}{lecturas}, $\theta
= 30/150 = 0.2$, contigs $= N e^{-6.4} = \num{267000}\times 0.00166
\approx 440$.
\end{solution}

\begin{solution}{exo:b3:genomics:6}
$P(X < 8) = P(X \le 7) \approx P\bigl(Z < (7.5 - 15)/2.74\bigr) =
P(Z < -2.74) \approx 0.003$. A $30\times$ los dos alelos de un
heterocigoto se ven casi siempre muchas veces, la cobertura es desigual
(las regiones ricas en GC reciben menos \omterm{def:b3:genomics:ngs}{lecturas}, de modo que algunos
sitios ven la mitad de la media) y quedan \omterm{def:b3:genomics:ngs}{lecturas} suficientes para
separar un alelo real de los errores de $10^{-3}$.
\end{solution}

\begin{solution}{exo:b3:genomics:7}
Toda \omterm{def:b3:genomics:ngs}{lectura} de \qty{150}{bp} del interior de la repetición es idéntica
venga de la copia que venga, de modo que el ensamblador construye un solo
nodo de \qty{6}{kb} con $500$ caminos de entrada y $500$ de salida; no
puede saber qué entrada se empareja con qué salida, y el \omterm{def:b3:genomics:assembly}{ensamblaje} se
rompe en $500$ huecos, con la repetición presente una sola vez y con una
cobertura $500$ veces la media. Las \omterm{def:b3:genomics:ngs}{lecturas} más largas que la repetición
más los flancos únicos de ambos lados --- \qty{8}{kb} o más --- abarcan
cada copia y la resuelven.
\end{solution}

\begin{solution}{exo:b3:genomics:8}
$4.8\times 10^{7}/1000 = \num{48000}$ diferencias codificantes; un
tercio, unas \num{16000}, alteran la proteína.
\end{solution}

\begin{solution}{exo:b3:genomics:9}
Falsos positivos esperados, $10^{6}\times 5\times 10^{-8} = 0.05$. Un
riesgo relativo de $1.15$ sobre una enfermedad del \qty{2}{\%} cambia
unos pocos casos por millar; mil personas contienen unos veinte casos,
demasiado pocos para verlo (la potencia crece con el número de casos y
con el cuadrado del efecto). Para un individuo, $0.3$ puntos porcentuales
no significan nada. Pero el alelo señala un gen cuyo cambio modesto de
actividad altera el riesgo de enfermedad --- el gen, y su vía, pueden ser
una diana farmacológica cuya inhibición completa tenga un efecto grande.
\end{solution}

\begin{solution}{exo:b3:genomics:10}
De genética de poblaciones: en las bacterias, con poblaciones enormes, la
selección elimina incluso las inserciones ligeramente costosas; en los
mamíferos, con poblaciones efectivas pequeñas, la deriva deja que se
acumule ADN no codificante levemente deletéreo --- lo apoya la
correlación inversa entre \omterm{prop:b3:genomics:sizes}{tamaño del genoma} y tamaño de población entre
linajes, y lo debilitan las excepciones. Estructural: los genomas
eucariotas los invaden transposones que las bacterias, con su sesgo hacia
la deleción y sin refugio meiótico para los elementos egoístas, depuran
--- lo apoya la correlación entre el \omterm{prop:b3:genomics:sizes}{tamaño del genoma} y el contenido de
transposones. Reguladora: un desarrollo complejo necesita más ADN
regulador --- lo apoyan los \omterm{def:b3:genomics:comparative}{elementos no codificantes conservados}
alrededor de los genes del desarrollo, y lo debilita el hecho de que solo
el \qty{8}{\%} del genoma humano muestre selección alguna, de modo que la
mayor parte del ADN no codificante no es regulador.
\end{solution}

\begin{solution}{exo:b3:genomics:11}
Los inicios de ambos tipos caen con densidad total $\rho = (N_{1} +
N_{2})/G$. Una \omterm{def:b3:genomics:ngs}{lectura} de longitud $L_{i}$ es la más a la derecha de su
contig si ninguna \omterm{def:b3:genomics:ngs}{lectura} de cualquier tipo empieza en las $L_{i} - T$
bases que la siguen: probabilidad $e^{-\rho(L_{i} - T)}$. De ahí que los
contigs sean $= N_{1}
e^{-\rho(L_{1} - T)} + N_{2} e^{-\rho(L_{2} - T)}$. Una \omterm{def:b3:genomics:ngs}{lectura} larga es
la más a la derecha con una probabilidad exponencialmente menor que una
corta, y cada \omterm{def:b3:genomics:ngs}{lectura} larga elimina además la posibilidad de un hueco a
lo largo de toda su longitud; el mismo número de bases en \omterm{def:b3:genomics:ngs}{lecturas} cortas
añade muchos puntos de inicio, pero cada ventana protegida es corta, de
modo que las \omterm{def:b3:genomics:ngs}{lecturas} largas ganan.
\end{solution}

\begin{solution}{exo:b3:genomics:12}
Dos duplicaciones del genoma completo predicen que el patrón cuádruple
sea de todo el genoma: cuartetos de segmentos cromosómicos \omterm{def:b3:genomics:comparative}{parálogos}
(paralogones) que llevan las mismas familias génicas en el mismo orden,
con las duplicaciones fechadas en el mismo momento para todas las
familias; un solo agrupamiento en el anfioxo y en \emph{Ciona}, que
divergieron antes de los sucesos; y en las lampreas, que divergieron
alrededor de ellos, un estado intermedio o derivado de forma
independiente. Las duplicaciones independientes solo del agrupamiento Hox
predicen paralogones únicamente para Hox, vecinos con historias distintas
y fechas de duplicación diferentes entre familias. Los paralogones de
todo el genoma y la datación compartida, tal como se observan, favorecen
la \omterm{def:b3:genomics:comparative}{duplicación del genoma completo}.
\end{solution}

\begin{solution}{pb:b3:genomics:1}
\textbf{1.} $0.002\times 6\times 10^{8} = 1.2\times 10^{6}$ \omterm{def:b3:genomics:ngs}{lecturas};
$c = 1.2\times 10^{6}\times 150/4.6\times 10^{6} \approx 39$.
\textbf{2.} $e^{-39} \approx 10^{-17}$: en la práctica no queda ninguna
base sin cubrir.
\textbf{3.} $\theta = 0.2$; contigs $= 1.2\times 10^{6} e^{-31} \approx
0$: un solo contig en teoría, con huecos solo en las repeticiones.
\textbf{4.} $c = \num{30000}\times 150/4.6\times 10^{6} = 0.98$; sin
cubrir, $e^{-0.98} = 0.38$, unas \qty{1.7}{Mb}; contigs $= \num{30000}
\times e^{-0.78} \approx \num{13700}$.
\textbf{5.} Los siete operones dan \omterm{def:b3:genomics:ngs}{lecturas} idénticas y colapsan en un
solo contig de \qty{5}{kb}, al que entran siete flancos izquierdos únicos
y del que salen siete flancos derechos cuyo emparejamiento se desconoce:
$14$ extremos de contig y siete huecos.
\textbf{6.} Unos \num{4600} genes; codificantes, $0.88\times 4.6\times 10^{6}
= 4.0\times 10^{6}$ bp, con un gen medio de \qty{880}{bp}.
\textbf{7.} $0.998\times 6\times 10^{8}\times 150/3.2\times 10^{9}
\approx 28\times$.
\textbf{8.} Unas $14$ \omterm{def:b3:genomics:ngs}{lecturas} por alelo.
\textbf{9.} $28\times 10^{-3}/3 \approx 0.009$ \omterm{def:b3:genomics:ngs}{lecturas} que muestran una
base equivocada dada --- la probabilidad de tres o más es de unos
$10^{-7}$ --- y el \qty{20}{\%} de $28$ son $5.6$ \omterm{def:b3:genomics:ngs}{lecturas}; un error no
pasa ninguna de las dos pruebas, mientras que de un alelo heterocigótico
real se esperan $14$.
\textbf{10.} $3.2\times 10^{9}/1500 \approx 2.1\times 10^{6}$ sitios
heterocigóticos; en secuencia codificante, $4.8\times 10^{7}/1500 \approx
\num{32000}$.
\textbf{11.} Una \omterm{def:b3:genomics:ngs}{lectura} de \qty{150}{bp} de un Alu casa casi igual de
bien con muchas de los $1.1$ millones de copias, de modo que se coloca en
la copia equivocada o se le da poca confianza de mapeo. Las diferencias
entre copias parecen entonces variantes heterocigóticas, y las variantes
verdaderas quedan escondidas entre ellas: las \omterm{met:b3:genomics:pipeline}{llamadas de variantes}
dentro de los Alu suelen filtrarse, y esos elementos son puntos ciegos de
la secuenciación de \omterm{def:b3:genomics:ngs}{lectura} corta.
\textbf{12.} $1.2\times 10^{-8}\times 6.4\times 10^{9} \approx 77$
mutaciones nuevas. Unas $14$ \omterm{def:b3:genomics:ngs}{lecturas} deberían mostrar la variante en el
niño; pero si el sitio de un progenitor solo está cubierto por cinco
\omterm{def:b3:genomics:ngs}{lecturas}, un alelo heterocigótico se pierde con probabilidad
$2^{-5} = 3\,\%$ y una variante heredada se llama por error como nueva
--- de ahí que los progenitores deban secuenciarse tan a fondo como el
niño.
\textbf{13.} $4.8\times 10^{7}\times 100/150 = 3.2\times 10^{7}$
\omterm{def:b3:genomics:ngs}{lecturas}, veinte veces menos que el genoma: el coste de secuenciación cae
veinte veces, más de lo que cuesta el paso de captura.
\textbf{14.} $1.1\times 10^{6}\times 300 = 3.3\times 10^{8}$ bp, cerca
del \qty{10}{\%} del genoma; el \qty{10}{\%} de las \omterm{def:b3:genomics:ngs}{lecturas}, unos $6\times
10^{7}$.
\textbf{15.} $\num{60000}\times 1500 = 9\times 10^{7}$ bp; $9\times
10^{7}/1.6\times 10^{10} = 0.56\,\%$.
\textbf{16.} $0.012\times 2.9\times 10^{9} = 3.5\times 10^{7}$
sustituciones, $1.7\times 10^{7}$ por linaje; tasa $1.7\times
10^{7}/(2.9\times 10^{9}\times 6.5\times 10^{6}) = 9\times 10^{-10}$
por base y año, $2.3\times 10^{-8}$ por generación de 25 años --- unas
dos veces la tasa de pedigrí de $1.2\times 10^{-8}$, una discrepancia
conocida que sugiere generaciones más largas en el pasado o una
separación más antigua.
\textbf{17.} Cuatro copias de \num{16000} serían \num{64000}; quedan
\num{20000}, de modo que de las \num{48000} copias adicionales solo
sobreviven \num{4000}: se perdió el \qty{92}{\%} de los duplicados.
\textbf{18.} El gusano y el ser humano tienen el mismo número de genes;
la complejidad está en cómo se usan --- el corte y empalme alternativo
multiplica un gen en miles de proteínas (\num{38016} para \emph{Dscam}),
la regulación combina factores de transcripción y potenciadores de
maneras propias de cada tipo celular, los \omterm{def:b3:rna-regulation:ncrna}{ARN no codificantes} añaden
capas y las proteínas interactúan en redes.
\textbf{19.} Elementos reguladores (potenciadores, promotores,
aislantes), genes de \omterm{def:b3:rna-regulation:ncrna}{ARN no codificante}, señales de corte y de
localización, orígenes y secuencias estructurales. Un potenciador
candidato se pone a prueba colocándolo delante de un gen indicador en un
embrión transgénico y mirando dónde se expresa el indicador, y después
suprimiendo el elemento del genoma con CRISPR y midiendo los genes
vecinos.
\textbf{20.} $0.05/10^{6} = 5\times 10^{-8}$; con $p < 10^{-5}$, $10^{6}
\times 10^{-5} = 10$ falsos positivos esperados.
\textbf{21.} No portador: $0.009/0.991 = 0.00908$ de posibilidades; con
una copia, $0.00954$, riesgo $0.945\,\%$; con dos copias, $0.01001$,
riesgo $0.99\,\%$.
\textbf{22.} La puntuación suma, sobre los $200$ alelos, el número de
copias de riesgo que lleva una persona ponderado por el logaritmo de la
razón de posibilidades de cada alelo. El recuento tiene media
$200\times 2\times 0.3 = 120$ y desviación típica
$\sqrt{200\times 2\times 0.3\times 0.7} \approx 9$; el \qty{1}{\%}
superior lleva unos $21$ alelos más que la media, y $1.05^{21} \approx
2.8$: casi el triple de posibilidades que una persona media, a partir de
alelos que individualmente cambian el riesgo un \qty{5}{\%}.
\textbf{23.} Los alelos de un bloque de decenas de kilobases se heredan
juntos (desequilibrio de ligamiento), de modo que cualquiera de ellos
muestra la misma asociación; el SNP genotipado es simplemente el que
estaba en la matriz. Para encontrar la variante causal: secuénciese la
región en casos y controles, redúzcase el conjunto por estadística y
pruébese después cada candidata --- ensayos con gen indicador para la
actividad potenciadora de cada alelo, edición de cada alelo en células y
medida de la expresión de los genes vecinos, colocalización con señales
de rasgos cuantitativos de expresión.
\textbf{24.} Se lee el $1 - e^{-0.5} = 39\,\%$ del genoma. Da una muestra
directa de una población en un momento conocido del pasado ---
frecuencias alélicas anteriores a migraciones y mestizajes posteriores,
la longitud de los segmentos neandertales (largos, porque pocas
generaciones de recombinación los habían roto) y con ello una fecha para
el mestizaje ---, cosa que los genomas modernos solo pueden inferir a
través de modelos.
\textbf{25.} Unos \num{13700} contigs en el \omterm{def:b3:genomics:assembly}{ensamblaje} de ensayo;
$28\times$ por genoma haploide, unas $14$ \omterm{def:b3:genomics:ngs}{lecturas} por alelo; y unas
$77$ mutaciones nuevas en un niño.
\end{solution}
